\section{Introduction}

Chest X-ray interpretation is one of the most commonly performed radiological tasks, yet it is resource-intensive and subject to variability among readers.
Developing reliable automated systems for detecting pulmonary and cardiac pathologies from chest X-ray images could support radiologists in clinical workflows, particularly in settings with limited specialist access.

The NIH ChestX-ray14 dataset~\cite{wang2017chestxray14} introduced a large-scale benchmark of 112,120 frontal chest X-ray images labelled with 14 disease categories using natural language processing applied to radiology reports.
A key challenge of this dataset is that it defines a \emph{multi-label} classification problem: each image may simultaneously present zero or more of the 14 pathologies.
This stands in contrast to single-label classification, since diseases frequently co-occur (e.g., Consolidation with Infiltration, or Effusion with Edema), and the model must predict each class independently.

A further challenge is \emph{class imbalance}.
Some pathologies are common—Infiltration appears in approximately 18\% of images—while others are rare, such as Hernia ($<$0.2\%).
Standard cross-entropy loss treats all samples equally, which can cause the model to neglect minority classes.
Loss function design choices, including asymmetric and focal variants, directly address this issue.

Recent work has continued to advance multi-label chest X-ray classification through novel loss functions and architectural improvements~\cite{hanif2025enhancing}.
Building on the CheXNet~\cite{rajpurkar2017chexnet} paradigm of using DenseNet121 pretrained on ImageNet as the backbone for this task, this work investigates four axes of variation:

\begin{enumerate}
  \item \textbf{Preprocessing}: the effect of CLAHE contrast enhancement.
  \item \textbf{Attention}: CBAM~\cite{woo2018cbam} module placement at different depths within the DenseNet121 feature hierarchy.
  \item \textbf{Loss function}: Binary Cross-Entropy (BCE), Asymmetric Loss (ASL)~\cite{ridnik2021asl}, and Focal Loss.
  \item \textbf{Optimizer}: Adam versus AdamW.
\end{enumerate}

The primary evaluation metric is Mean Test AUC (macro-average ROC-AUC over all 14 classes), which is robust to class imbalance and threshold-independent.
Precision, Recall, and F1 score at threshold 0.5 are reported as secondary metrics to characterize sensitivity at the operating threshold.

This paper reports results from ten experimental configurations, identifies the best-performing models under both AUC-oriented and Recall/F1-oriented objectives, and provides qualitative interpretability analysis via Grad-CAM.
